\chapter*{Preface: mathematics, machines and exposition}
\addcontentsline{toc}{chapter}{Preface: mathematics, machines and exposition}
\label{tx:preface:provenance}
\begingroup\fontsize{10.5}{12.6}\selectfont
AI is changing how mathematics can be explored, checked and explained.
This project suggests a corresponding way to present mathematical research:
an interactive paper, with further explanations attached to its arguments,
and a tutorial textbook which develops the ideas at a human pace.
The two formats serve different purposes. The paper states the results
and their proofs concisely; the textbook supplies the definitions,
examples and intermediate reasoning needed to understand the whole route.

Machine-generated proofs can be hard to follow. They may combine ideas
from several research communities without explaining why a construction
was chosen or how its pieces fit together. The initial proofs generated
by Astra and Claude also omitted several key references to prior human
work. Their mathematical knowledge reflects a body of human research,
although the outputs do not establish which particular sources were
used in training or how they influenced the constructions. A formally
checked argument therefore also needs a readable account of the
mathematics and its antecedents. This book provides that account for
the linear degree gap.

This is an \emph{AI-generated tutorial textbook}, developed under my
mathematical direction, supervision and editing. Astra supplied much of
the drafting and restructuring. Astra and Claude also contributed
substantially to the proof search and formalisation. I selected the
programme, supplied my earlier work and mathematical input, challenged
arguments, read successive chapters and requested clearer explanations
and worked examples. This continuing exchange shaped both the proof
development and its presentation.

The published mathematical starting point is my LICS~2008
paper~\cite{Riis}; the 53-page 2009 report supplies a longer account
\cite{RiisLong}. The conference paper established a four-way
classification, or \emph{quadrichotomy}, of uniform algebraic
refutation complexity: $\mathrm T_1,\ldots,\mathrm T_4$, with
periodicity, NS--PC type agreement and undecidability for every
nontrivial union of types, including the no-finite-model promise.
The logarithmic and double-logarithmic scales used
Kraj\'i\v cek's Theorem~3.5~\cite{Krajicek}, based on Ajtai's
symmetric-system work~\cite{Ajtai}.
My joint journal papers with Meera Sitharam develop uniform
translations~\cite{RiisSitharamTautologies} and uniformly generated
permutation submodules over fields of characteristic
zero~\cite{RiisSitharamModules}. A separate 1997 announcement proposed
a nonconstant-to-linear degree gap. Its inducing-down argument has
an unresolved gap, as I recorded in the 2009 report
\cite[p.~50, footnote~3; p.~53, reference~33]{RiisLong}; this concerns
the proof and does not refute the asserted theorem. My later
unpublished notes pursued stronger bounds and suggested lifting
separating vectors between ranks, with important steps incomplete.
The aim here was to obtain complete proofs of linear degree bounds in
the universe size, in every fixed positive characteristic, and to treat
uniform polynomial consequences as well as refutations.

The central proof search included, in my recollection, roughly two to
three days of concentrated work with Astra and Claude, alongside my
interactions. Checking, repairs, extensions and exposition continued
afterwards. When the construction reached an $\Omega(\sqrt n)$ bound in
characteristic two, I pressed for an $\Omega(n)$ bound and an argument
covering all positive characteristics. The difficult work was to build
one lifting mechanism compatible with all the required operations.

\newpage
The development and subsequent comparison with the literature brought
together proof complexity, model theory, modular representation theory
and algebraic combinatorics. Some of these connections were unfamiliar
to me; the assistants helped bring them into the discussion.
Kraj\'i\v cek's earlier work already treated uniform polynomial targets,
Euler structures and bounded-degree NS--PC comparison, including the
NS-span plateau argument~\cite{KrajEuler}\cite[Theorem~5.5]{Krajicek}. The wider context
includes Church--Ellenberg--Farb's FI framework for representations as
the number of labels grows, Harman's modular periodicity in a linear
range, and Hone and coauthors' explicit periodic functors through
semisimplification~\cite{CEF,Harman,HKPSYZ}. Harman's equivalences
already respect induction, restriction, tensor products and
duality~\cite[Propositions~2.7 and~2.12]{Harman}. The construction
here proves the required closure preservation and target detection
on the actual polynomial maps with its own quantitative bounds.
Integral centraliser and
rook combinatorics provide another connection~\cite{BDM}: realising
compatible linear maps as combinations of label permutations. These
communities study stability as the number of labels grows and the
construction of maps between representations.
Chapter~\ref{tx:chap:contract} explains these antecedents; the finite
constructions required for our polynomial proof rules are proved here.

The principal construction developed in this project is in
Chapters~\ref{tx:ch:linear-layers}--\ref{tx:ch:prime}: compatible
corrections reuse one retained pool of labels, and the odd-characteristic
argument completes the same linear budget.
Chapter~\ref{tx:chap:assemble-transfer} transports the actual axiom,
multiplication and indexed-target maps.
Chapter~\ref{tx:sec:comparison} independently supplies the comparison's
uniform chain bound. The four qualitative types are inherited from my
earlier work; the present quantitative bounds are linear in
$\mathrm T_2$, linear on the hard residues in $\mathrm T_3$, and give
finite linear peaks in $\mathrm T_4$.
Chapter~\ref{tx:sec:undecidability} gives explicit reductions separating
any two types under the no-finite-model promise already stated in
LICS~2008, refining its every-union formulation. This treatment was
developed interactively by me and Astra from my earlier undecidability
arguments. The NS proof system is due to Beame, Impagliazzo,
Kraj\'i\v cek, Pitassi and Pudl\'ak~\cite{NSFoundations}; PC is due
to Clegg, Edmonds and Impagliazzo~\cite{PCFoundations}.
The induction-chain separation used in
Chapter~\ref{tx:ch:running-order} is due to Buss and
Pitassi~\cite{BussPitassiInduction}; the design method is developed
in Buss~\cite{BussDesign}, and the modular bijection antecedent is
Beame--Riis~\cite{BeameRiis}. The compression discussion belongs to
separate companion work.

I developed this tutorial interactively by reading drafts, identifying
steps I could not follow, proposing abstractions and requesting examples
that would carry the reader through the general argument. Alongside
supplying my earlier manuscripts and guiding the proof search, an
important part of my contribution was to identify relevant prior human
work, check the connections against the literature, and ensure that
these contributions received proper credit. My intention is to make
the textbook available on GitHub alongside the interactive
version of the paper, and potentially use it for lectures and
presentations. The contribution to exposition is this coordinated
combination: a scientific paper with expandable explanations and a
longer, progressively developed tutorial.

I stated and proved the following undecidability theorem: for every
nonempty proper subset $A\subsetneq\{1,2,3,4\}$, it is undecidable
whether a uniform family arising from a first-order sentence belongs
to one of the types $\mathrm T_i$ with $i\in A$.
This theorem appears in my LICS~2008 paper, with further details in
the 2009 report~\cite{Riis,RiisLong}.
The no-finite-model promise is explicitly part of its published
statement~\cite[Theorem~1]{Riis}. Its Lemma~8 gives a
$\mathrm T_1$-versus-prescribed-type switch; the controlled-join
construction in Chapter~\ref{tx:sec:undecidability} supplies the
remaining pairwise reductions. This comparison concerns that
published formulation and inspected construction.
The core degree, comparison and undecidability results now have Lean
formalisations and independent source, statement and axiom audits.
Newly added teaching examples and the topological reformulation are
written deductions; they are not additional Lean theorem endpoints. The book explains
the mathematics and its provenance, while making explicit how its text
was generated and how I directed and revised it.

\par\endgroup
